\documentclass[11pt]{article}
\usepackage[T1]{fontenc}
\usepackage{lmodern}
\usepackage{microtype}
\usepackage{amsmath, amssymb}
\usepackage{geometry}
\geometry{margin=1in}
\usepackage{hyperref}
\hypersetup{colorlinks=true, linkcolor=blue, urlcolor=blue, citecolor=blue}

\title{Input-Anchored Cost Model for R/A/B/C Hamiltonian Application}
\author{Uni20 Developer Documentation}
\date{\today}

\begin{document}

\maketitle

\section{Purpose}

This note defines a restricted, implementation-oriented cost model for applying
the two-site DMRG effective Hamiltonian in R/A/B/C form.  It is a simplification
of the more general hypergraph placement model in
\texttt{docs/latex/rabc\_hypergraph\_partitioning.tex}.  The aim is to remove
most underdetermined communication choices while keeping the dominant
Hamiltonian scheduling decisions.

The effective Hamiltonian is represented by sparse terms
\begin{equation}
  t = (r(t),a(t),b(t),c(t),f_t),
  \qquad
  R_{r(t)} \mathrel{+}= f_t A_{a(t)} B_{b(t)} C_{c(t)} .
\end{equation}
The Hamiltonian Lanczos path has one important invariant: the input center
vector $B$ and output center vector $R$ have the same block space and must use
the same persistent layout, so the next matvec can consume the output without a
relayout.

\section{Anchoring Assumption}

Let $\mathcal{D}$ be the device set and $\mathcal{K}$ be the center-block set.
Choose one canonical owner
\begin{equation}
  h : \mathcal{K} \to \mathcal{D}.
\end{equation}
For Hamiltonian application, $B_k$ is initially resident on $h(k)$ and $R_k$
must finally be resident on $h(k)$.

The input-anchored rule is:
\begin{equation}
  \text{the first GEMM for term } t \text{ runs on } h(b(t)).
\end{equation}
The output rule is:
\begin{equation}
  \text{the second GEMM for term } t \text{ runs on } h(r(t)).
\end{equation}
These two rules remove the independent choice of where to compute first-stage
products, where to assemble temporaries, and where to accumulate final outputs.
Given a layout $h$ and a left/right path choice for each term, the copy and
temporary schedule is induced.

This is not the fully general communication schedule.  It excludes moving the
first GEMM away from the input block, moving the second GEMM away from the final
output block, tree reductions, temporary recomputation instead of copying, and
arbitrary environment replication.  These are deliberate restrictions.  They
turn an ill-determined communication problem into a concrete cost functional.

\section{Optimization Variables}

For each term choose only its contraction side:
\begin{equation}
  s_t \in \{L,R\}.
\end{equation}
Here $s_t=L$ means left-first,
\begin{equation}
  A_aB_b \quad \text{then} \quad (A_aB_b)C_c,
\end{equation}
and $s_t=R$ means right-first,
\begin{equation}
  B_bC_c \quad \text{then} \quad A_a(B_bC_c).
\end{equation}
The full restricted decision is therefore
\begin{equation}
  (h,\{s_t\}_{t\in\mathcal{T}}).
\end{equation}

\section{Induced DAG}

The induced DAG uses conservative exact temporaries.  The scalar coefficient
$f_t$ is applied after the first-stage GEMM product has been formed, while
accumulating that product into the local temporary.
Before the resident matvec begins, setup has already materialized every
environment block needed on the devices where the GEMM kernels below will run.

\subsection{Left-First Terms}

For each device, first form the distinct left-first GEMM products needed by
terms whose input block lives on that device:
\begin{equation}
  G_L(d) =
    \{(a(t),b(t)) : s_t=L,\ h(b(t))=d\}.
\end{equation}
For every $(a,b)\in G_L(d)$, compute once on device $d$:
\begin{equation}
  X^d_{a,b} = A_a B_b .
\end{equation}
The block $B_b$ is already on $d=h(b)$.  The environment block $A_a$ is made
available on $d$ by the setup procedure before the resident matvec is timed.

Then, for every left-first term $t$ with
$(r,a,b,c)=(r(t),a(t),b(t),c(t))$, accumulate on $h(b)$:
\begin{equation}
  Q^{L,h(b)}_{r,c}
  \mathrel{+}= f_t X^{h(b)}_{a,b}.
\end{equation}

For each nonzero local partial $Q^{L,d}_{r,c}$, copy it to $h(r)$ if
$d\ne h(r)$ and accumulate it into the assembled temporary
\begin{equation}
  Q^L_{r,c} = \sum_{d\in\mathcal{D}} Q^{L,d}_{r,c}
  \quad \text{on } h(r).
\end{equation}
Then execute the second GEMM on $h(r)$:
\begin{equation}
  R_r \mathrel{+}= Q^L_{r,c} C_c
  \quad \text{on } h(r).
\end{equation}
The environment block $C_c$ has also been made available on $h(r)$ by setup.

\subsection{Right-First Terms}

For each device, first form the distinct right-first GEMM products needed by
terms whose input block lives on that device:
\begin{equation}
  G_R(d) =
    \{(b(t),c(t)) : s_t=R,\ h(b(t))=d\}.
\end{equation}
For every $(b,c)\in G_R(d)$, compute once on device $d$:
\begin{equation}
  Y^d_{b,c} = B_b C_c .
\end{equation}
The block $B_b$ is already on $d=h(b)$.  The environment block $C_c$ is made
available on $d$ by the setup procedure before the resident matvec is timed.

Then, for every right-first term $t$ with
$(r,a,b,c)=(r(t),a(t),b(t),c(t))$, accumulate on $h(b)$:
\begin{equation}
  Q^{R,h(b)}_{r,a}
  \mathrel{+}= f_t Y^{h(b)}_{b,c}.
\end{equation}

For each nonzero local partial $Q^{R,d}_{r,a}$, copy it to $h(r)$ if
$d\ne h(r)$ and accumulate it into the assembled temporary
\begin{equation}
  Q^R_{r,a} = \sum_{d\in\mathcal{D}} Q^{R,d}_{r,a}
  \quad \text{on } h(r).
\end{equation}
Then execute the second GEMM on $h(r)$:
\begin{equation}
  R_r \mathrel{+}= A_a Q^R_{r,a}
  \quad \text{on } h(r).
\end{equation}
The environment block $A_a$ has also been made available on $h(r)$ by setup.

\section{Cost Functional}

The objective is predicted makespan:
\begin{equation}
  \min_{h,s} \max_{d\in\mathcal{D}} C_d(h,s).
\end{equation}
For each device $d$,
\begin{equation}
  C_d =
    C^{\mathrm{gemm}}_d
  + C^{\mathrm{accum}}_d
  + C^{\mathrm{temp}}_d
  + C^{\mathrm{overhead}}_d
  + C^{\mathrm{memory}}_d.
\end{equation}

\subsection{GEMM Cost}

Let $\Phi(A_a,B_b)$ denote the GEMM flop count for $A_aB_b$ and similarly for
other products.  The first-stage cost is
\begin{align}
  C^{\mathrm{first}}_d &=
    \alpha_{L,1}
      \sum_{(a,b)\in G_L(d)}
        \Phi(A_a,B_b) \notag\\
  &\quad+
    \alpha_{R,1}
      \sum_{(b,c)\in G_R(d)}
        \Phi(B_b,C_c).
\end{align}
The second-stage cost is grouped by assembled temporaries:
\begin{align}
  C^{\mathrm{second}}_d &=
    \alpha_{L,2}
      \sum_{\substack{(r,c): Q^L_{r,c}\ne 0\\h(r)=d}}
        \Phi(Q^L_{r,c},C_c) \notag\\
  &\quad+
    \alpha_{R,2}
      \sum_{\substack{(r,a): Q^R_{r,a}\ne 0\\h(r)=d}}
        \Phi(A_a,Q^R_{r,a}).
\end{align}
Then
\begin{equation}
  C^{\mathrm{gemm}}_d =
    C^{\mathrm{first}}_d + C^{\mathrm{second}}_d .
\end{equation}
The coefficients $\alpha$ should be fitted, not assumed from peak FLOPS, because
small block GEMMs include cuBLAS dispatch and launch overhead.

\subsection{Temporary Accumulation Cost}

After forming each first-stage product, the scheduler scales it by $f_t$ and
accumulates it into a local temporary.  This is lower order than the GEMM, but
it can matter for many small symmetry blocks.  Let
$\operatorname{bytes}(X^d_{a,b})$ and $\operatorname{bytes}(Y^d_{b,c})$ be the
temporary product sizes.  A simple bandwidth-like model is
\begin{align}
  C^{\mathrm{accum}}_d &=
    \delta_L
      \sum_{\substack{t\in\mathcal{T}\\s_t=L\\h(b(t))=d}}
        \operatorname{bytes}(X^d_{a(t),b(t)}) \notag\\
  &\quad+
    \delta_R
      \sum_{\substack{t\in\mathcal{T}\\s_t=R\\h(b(t))=d}}
        \operatorname{bytes}(Y^d_{b(t),c(t)}).
\end{align}
The associated kernel-launch or loop overhead is represented separately by the
count features below.

\subsection{Environment Setup}

The first GEMM needs one environment side on $h(b)$; the second GEMM needs the
other side on $h(r)$.  In this input-anchored resident matvec model, environment
placement is part of setup, not part of the Hamiltonian-apply cost functional.
The setup step must materialize the required distinct environment blocks:
\begin{align}
  U^{L,1}_A(d) &= \{a(t): s_t=L,\ h(b(t))=d\}
    &&\text{for } A_aB_b \text{ products},\\
  U^{L,2}_C(d) &= \{c(t): s_t=L,\ h(r(t))=d\}
    &&\text{for } Q^L_{r,c}C_c \text{ products},\\
  U^{R,1}_C(d) &= \{c(t): s_t=R,\ h(b(t))=d\}
    &&\text{for } B_bC_c \text{ products},\\
  U^{R,2}_A(d) &= \{a(t): s_t=R,\ h(r(t))=d\}
    &&\text{for } A_aQ^R_{r,a} \text{ products}.
\end{align}
The setup benchmark may separately measure bytes and latency for this
materialization, but those costs should not be included in $C_d$ when fitting
or optimizing the steady-state Lanczos matvec.

\subsection{Temporary Migration}

Left-first temporaries are keyed by $(r,c)$ and initially accumulated on every
device that owns at least one contributing $B_b$.  Define
\begin{align}
  S^L_{r,c} &= \{h(b(t)): s_t=L,\ r(t)=r,\ c(t)=c\},\\
  S^R_{r,a} &= \{h(b(t)): s_t=R,\ r(t)=r,\ a(t)=a\}.
\end{align}
The temporary migration cost charged to destination device $d$ is
\begin{align}
  C^{\mathrm{temp}}_d &=
    \beta_P
    \sum_{\substack{(r,c): Q^L_{r,c}\ne 0\\h(r)=d}}
      \sum_{\substack{s\in S^L_{r,c}\\s\ne d}}
        \operatorname{bytes}(Q^L_{r,c}) \notag\\
  &\quad+
    \beta_P
    \sum_{\substack{(r,a): Q^R_{r,a}\ne 0\\h(r)=d}}
      \sum_{\substack{s\in S^R_{r,a}\\s\ne d}}
        \operatorname{bytes}(Q^R_{r,a}).
\end{align}
This model assumes local accumulation before transfer.  That is not a proof of
optimality, but it avoids shipping one temporary per term and makes the
communication cost quadratic-order in block dimensions rather than cubic-order
GEMM work.

\subsection{CUDA and Scheduler Overheads}

For small blocks, count features are necessary.  A minimal overhead model is
\begin{equation}
  C^{\mathrm{overhead}}_d =
    \gamma_1 N^{\mathrm{first\_gemm}}_d
  + \gamma_A N^{\mathrm{term\_accum}}_d
  + \gamma_2 N^{\mathrm{second}}_d
  + \gamma_P N^{\mathrm{temp\_copy}}_d
  + \gamma_Z N^{\mathrm{zero}}_d .
\end{equation}
Here $N^{\mathrm{first\_gemm}}_d=|G_L(d)|+|G_R(d)|$ is the number of cached
first-stage GEMM products on $d$, $N^{\mathrm{term\_accum}}_d$ is the number
of term contributions accumulated into local temporaries on $d$,
$N^{\mathrm{second}}_d$ is the number of assembled second-stage temporaries
multiplied on $d$, $N^{\mathrm{temp\_copy}}_d$ counts temporary migrations
into $d$, and $N^{\mathrm{zero}}_d$ counts output or temporary buffers that
must be initialized.  These coefficients should be learned from no-trace
fixture benchmarks or low-overhead CUDA event timing.

\subsection{Memory Pressure}

Let $M_d(h,s)$ be the peak live bytes for the resident environment working set,
local partial temporaries, assembled temporaries, and output buffers on $d$.
The production model should enforce
\begin{equation}
  M_d(h,s) \le M^{\max}_d.
\end{equation}
A prototype model may instead add
\begin{equation}
  C^{\mathrm{memory}}_d = \mu M_d(h,s),
\end{equation}
but hard capacity constraints are the correct final form.

\section{Layout Domain and Storage Packing}

This model does not require a one-dimensional block ordering.  The layout
$h(k)$ is a graph placement of logical center blocks, and the cost functional
above applies directly to arbitrary placements.  That is the intended target
for local search, graph partitioning, or offline integer programming.

Storage packing is a separate implementation step.  After $h$ is chosen, the
storage layer may pack all blocks assigned to a device into one or more
contiguous slabs.  Contiguity is therefore a storage consequence, not part of
the mathematical optimization domain.

\section{Environment-Construction Variant}

For environment construction, the input and output families need not share one
layout.  Use separate maps
\begin{equation}
  h_B : \mathcal{B} \to \mathcal{D},
  \qquad
  h_R : \mathcal{R} \to \mathcal{D}.
\end{equation}
The same input-anchored rule can run first GEMMs on $h_B(b)$ and second GEMMs
on $h_R(r)$.  The Hamiltonian constraint $h_B=h_R$ should not be imposed in
that case.

\end{document}
